% *==================================================================================*
% *                     Review vs. Camera-Ready settings                             *
% *==================================================================================*
%double-blind
\documentclass[cameraready]{Interspeech}

%==================================================================================
\title{Deep Conditioning via Timestep Embedding Injection\\for Multi-Degradation Speech Enhancement}

%==================================================================================
% orcid 가입하셔서
\author[affiliation={1}, orcid=0009-0002-3994-7018]{Seokhoon}{Moon}
\author[affiliation={2}, orcid=0009-0006-4946-6305]{Kyudan}{Jung}
\author[affiliation={2}, orcid=0000-0003-1071-4835, correspondingauthor]{Jaegul}{Choo}

%==================================================================================
\address{
    $^1$ KAIST
    $^2$ KAIST AI
}

%==================================================================================
\email{moonx010@kaist.ac.kr, kyudan@kaist.ac.kr, jchoo@kaist.ac.kr}

%==================================================================================
\keywords{Speech Enhancement, Diffusion Models, Multi-Degradation, Deep Conditioning, Timestep Embedding}

\usepackage{graphicx}
\usepackage{amsmath, amssymb, amsfonts, bm}
\usepackage{booktabs}
\usepackage{tikz}
\usetikzlibrary{positioning, arrows.meta, calc, fit, backgrounds}
\usepackage{multirow}
\usepackage{dblfloatfix}
\usepackage{stfloats}

%==================================================================================
\begin{document}

\maketitle

%==================================================================================
\begin{abstract}
Real-world speech signals are often corrupted by multiple degradation types simultaneously, including additive noise, reverberation, and nonlinear distortion.
While diffusion-based speech enhancement methods have achieved strong performance on single-degradation scenarios, they struggle with compound degradations.
Recent noise-aware approaches such as NASE inject noise embeddings via shallow input-level addition, which limits the model's ability to leverage degradation information throughout its deep architecture.
We investigate the effect of \textit{conditioning injection depth} in score-based speech enhancement and find that it is critical: in controlled experiments where only the injection method varies, shallow input addition \textit{degrades} ESTOI from 0.77 to 0.73---worse than using no encoder at all---while injecting the same conditioning into the \textit{timestep embedding}, which propagates through all ${\sim}37$ residual blocks, improves ESTOI to 0.80.
Building on this finding, we propose a multi-degradation framework with a WavLM encoder that jointly characterizes noise type, reverberation, and distortion via multi-task learning.
On compound degradations, the proposed method achieves ESTOI 0.80 and UTMOS 3.71, substantially outperforming both single-degradation baselines and the same model with input-level addition.
\end{abstract}


%==================================================================================
\section{Introduction}
%==================================================================================

Speech enhancement aims to recover clean speech from degraded observations.
While most existing methods focus on additive noise removal~\cite{fu2021metricgan, lu2023mpsenet, richter2023speech}, real-world speech is typically corrupted by a combination of degradations: additive noise from the environment, reverberation from room acoustics, and nonlinear distortion from recording equipment~\cite{su2020hifi, liu2022voicefixer}.
A unified model that handles all degradation types simultaneously is highly desirable but remains challenging.

Diffusion-based generative models have achieved strong results for speech enhancement~\cite{richter2023speech, lu2022conditional}.
SGMSE+~\cite{richter2023speech} models the conditional distribution of clean speech using score-based stochastic differential equations (SDEs) in the complex STFT domain.
To improve generalization, noise-aware methods inject external noise information as a conditioning signal.
NASE~\cite{hu2023nase} uses a BEATs~\cite{chen2023beats} encoder with noise classification (NC) loss to produce acoustic embeddings, while NADiffuSE~\cite{wang2023nadiffuse} extracts noise representations as global conditions.
More recently, DiTSE~\cite{zheng2025ditse} employs a Diffusion Transformer with Classifier-Free Guidance (CFG)~\cite{ho2022classifier}.

However, these methods share two critical limitations.
First, they only condition on additive noise, ignoring reverberation and distortion.
Second, they rely on \textit{shallow conditioning injection}: NASE adds embeddings at the input layer (a single injection point), which limits the influence of conditioning on the deep score network.
Our ablation experiments reveal that this shallow injection can be \textit{worse than no injection at all}: input-level conditioning yields lower performance than a model without any encoder on compound degradations (Section~\ref{sec:ablation}).

We present three contributions:
\begin{itemize}
\item \textbf{Empirical finding on injection depth:} Through controlled experiments where only the injection method varies (same data, same encoder), we show that shallow input addition---used by NASE~\cite{hu2023nase} and NADiffuSE~\cite{wang2023nadiffuse}---is \textit{worse than no conditioning at all} on compound degradations (ESTOI 0.73 vs.\ 0.77).
This challenges the implicit assumption that any form of conditioning injection is beneficial.

\item \textbf{Deep conditioning via timestep embedding injection:} We inject degradation conditioning into the \textit{timestep embedding}, which propagates through all ${\sim}37$ residual blocks of the NCSN++~\cite{song2020score} backbone.
Unlike FiLM~\cite{perez2018film} or AdaLN-Zero~\cite{peebles2023scalable}, which require per-layer modifications, our method reuses the backbone's existing conditioning pathway with no architectural changes.

\item \textbf{Multi-degradation conditioning framework:} We extend noise-aware SE from additive noise only to compound degradations.
A WavLM~\cite{chen2022wavlm} encoder with three specialized heads jointly characterizes noise type (96.7\% accuracy), reverberation level ($T_{60}$ correlation 0.98), and distortion intensity, supervised by multi-task auxiliary losses~\cite{caruana1997multitask}.
This provides explicit, interpretable degradation awareness and enables principled data augmentation combining noise, reverberation, and distortion.
\end{itemize}

Controlled experiments on VoiceBank-DEMAND~\cite{valentini2016investigating} with multi-degradation data demonstrate that deep injection is the decisive factor: adding an encoder with input addition degrades performance, while the same encoder with deep injection yields the best results across all metrics.


%==================================================================================
\section{Proposed Method}
%==================================================================================

% ---- Figure 1: Architecture ----
\begin{figure*}[!t]
\centering
\includegraphics[width=\textwidth]{figure1.png}
\caption{Overview of the proposed framework. A frozen WavLM encoder extracts features from the degraded speech, which are projected into three degradation-specific branches (noise, reverberation, distortion). The concatenated projections are mapped to a 512-dimensional vector and added to the timestep embedding (\textbf{temb}), which is injected into all ${\sim}37$ residual blocks of the NCSN++ backbone. Dashed arrows indicate auxiliary losses used only during training.}
\label{fig:architecture}
\end{figure*}


\subsection{Preliminaries: Score-Based Speech Enhancement}

We build upon SGMSE+~\cite{richter2023speech}, which models the clean speech distribution conditioned on the noisy observation using a score-based SDE framework.
The forward process gradually corrupts the clean signal $\mathbf{x}$ toward the noisy observation $\mathbf{y}$:
\begin{equation}
\mathrm{d}\mathbf{x}_t = \gamma(\mathbf{y} - \mathbf{x}_t)\,\mathrm{d}t + g(t)\,\mathrm{d}\mathbf{w}_t
\end{equation}
where $\gamma$ controls mean reversion and $g(t)$ is the diffusion coefficient.
A score network $\mathbf{s}_\theta(\mathbf{x}_t, \mathbf{y}, t) \approx \nabla_{\mathbf{x}_t} \log p_t(\mathbf{x}_t | \mathbf{y})$ is trained via denoising score matching, and clean speech is recovered by solving the reverse SDE with a predictor-corrector sampler.
The NCSN++~\cite{song2020score} backbone computes a timestep embedding $\mathbf{e}_t = \text{MLP}_\text{time}(t) \in \mathbb{R}^{512}$ that is added to features in every residual block, providing timestep awareness at all network depths.


\subsection{Multi-Degradation Encoder}

To characterize diverse degradation types, we employ a frozen WavLM Base~\cite{chen2022wavlm} encoder followed by a convolutional post-processing network that maps the 768-dimensional features to a 256-dimensional shared representation $\mathbf{h} \in \mathbb{R}^{256}$.

Three specialized heads operate on the WavLM features:
\begin{itemize}
\item \textbf{Noise head}: 11-class classification (10 DEMAND~\cite{thiemann2013demand} noise types + ``none'')
\item \textbf{Reverb head}: Regression predicting $T_{60}$ (reverberation time)
\item \textbf{Distortion head}: Regression predicting distortion intensity
\end{itemize}
These heads serve dual purposes: (1) providing auxiliary supervision to learn discriminative features, and (2) enabling interpretable per-degradation analysis at inference time.


\subsection{Deep Conditioning via Timestep Embedding Injection}

We exploit the NCSN++ backbone's existing timestep embedding mechanism for degradation conditioning.
The shared representation $\mathbf{h}$ is projected into three 128-dimensional branch-specific embeddings, concatenated, and mapped to match the timestep embedding dimension:
\begin{align}
\mathbf{c}_k &= \mathbf{W}_k \mathbf{h}, \quad k \in \{\text{noise, reverb, distort}\} \label{eq:branch} \\
\mathbf{c}_\text{extra} &= \text{MLP}_\text{cond}([\mathbf{c}_\text{noise}; \mathbf{c}_\text{reverb}; \mathbf{c}_\text{distort}]) \label{eq:mlp}
\end{align}
where $[\cdot;\cdot]$ denotes concatenation and $\text{MLP}_\text{cond}: \mathbb{R}^{384} \to \mathbb{R}^{512}$.
The conditioning is then injected by simple addition:
\begin{equation}
\tilde{\mathbf{e}}_t = \mathbf{e}_t + \mathbf{c}_\text{extra}
\label{eq:injection}
\end{equation}

This replaces NASE's input-level addition (1 injection point) with propagation through all ${\sim}37$ residual blocks, making every layer of the score network degradation-aware.
Crucially, this requires \textit{no architectural changes} to the backbone---only a single addition to the existing timestep embedding.


\subsection{Training Objective}

The total loss combines the score matching objective with auxiliary multi-task losses:
\begin{equation}
\mathcal{L} = \mathcal{L}_\text{score} + \lambda \left( \mathcal{L}_\text{noise} + \mathcal{L}_\text{reverb} + \mathcal{L}_\text{distort} \right)
\label{eq:loss}
\end{equation}
where $\mathcal{L}_\text{noise}$ is cross-entropy for noise classification, $\mathcal{L}_\text{reverb}$ and $\mathcal{L}_\text{distort}$ are MSE for regression, and $\lambda = 0.3$.

Following CFG~\cite{ho2022classifier}, each branch embedding is independently dropped to zero with probability $p = 0.1$ during training.
This enables the model to handle missing degradation types gracefully at inference, as it has learned to enhance speech with any subset of conditioning branches active.


%==================================================================================
\section{Experiments}
%==================================================================================

\subsection{Experimental Setup}

\subsubsection{Datasets}
We use VoiceBank-DEMAND~\cite{valentini2016investigating, thiemann2013demand} (16\,kHz) as the base dataset containing 11,572 training utterances with 10 noise types at SNRs of 0, 5, 10, and 15\,dB.
To create multi-degradation training data, we augment each clean utterance with combinations of: (i) additive noise from DEMAND, (ii) reverberation via synthetic room impulse responses ($T_{60} \in [0.3, 1.0]$\,s) generated with pyroomacoustics~\cite{scheibler2018pyroomacoustics}, and (iii) nonlinear distortion via soft clipping $f(x) = \tanh(\alpha x)/\tanh(\alpha)$ with intensity $\alpha \in [1.5, 5.0]$.
The augmented dataset contains 34,716 utterances spanning six degradation categories.
Although soft clipping is a controlled simulation, the diffusion-based generative process does not learn a simple mathematical inverse; instead, conditioned on the encoder's severity estimation, it reconstructs missing spectral components stochastically, allowing generalization to real-world distortions as demonstrated in the wild evaluation.

We evaluate on: (1) \textbf{Noise-only}: VoiceBank-DEMAND test (824 files), (2) \textbf{Multi-degradation}: 2,472 files covering all six categories, and (3) \textbf{Wild}: three real-world datasets---VOiCES (1,600 files, room recordings with noise and reverberation), DAPS (1,500 files, real device recordings), and URGENT (1,000 files, diverse degradations).

\subsubsection{Baselines}
We compare against five baselines trained on noise-only data: SGMSE+ (denoise)~\cite{richter2023speech}, SGMSE+ (dereverb)~\cite{richter2023speech}, MetricGAN+~\cite{fu2021metricgan}, MP-SENet~\cite{lu2023mpsenet}, and NASE~\cite{hu2023nase} (BEATs encoder with input-level addition).
To isolate the effect of the encoder and injection method, we also train two additional models on the same multi-degradation data: (1) SGMSE+ without any encoder (``SGMSE+ retrained''), and (2) our full model but with NASE-style input addition instead of deep injection.
These three models (no encoder, input addition, deep injection) share identical training data, enabling a controlled comparison.

\subsubsection{Metrics and Implementation}
We report PESQ~\cite{rix2001perceptual}, ESTOI~\cite{jensen2016algorithm}, SI-SDR~\cite{le2019sdr}, and UTMOS~\cite{saeki2022utmos} (a neural MOS predictor).
We use the NCSN++~\cite{song2020score} backbone with the OUVE SDE from SGMSE+~\cite{richter2023speech}.
Models are trained for 160 epochs with Adam ($\text{lr}=10^{-4}$), EMA decay 0.999, batch size 32 (8 GPUs $\times$ 4), and 30 reverse diffusion steps at inference.
The WavLM encoder is frozen throughout training.


\subsection{Main Results}

\begin{table*}[!t]
\centering
\caption{Speech enhancement results on noise-only and multi-degradation conditions. Models above the line are trained on noise-only data; models below the line use the same multi-degradation data, isolating the effect of the encoder and injection method. Best per condition in \textbf{bold}. SDR = SI-SDR (dB).}
\label{tab:main_results}
\small
\setlength{\tabcolsep}{4pt}
\begin{tabular}{l|cccc|cccc}
\toprule
& \multicolumn{4}{c|}{Noise-Only (824 files)} & \multicolumn{4}{c}{Multi-Degradation (2472 files)} \\
Method & PESQ & ESTOI & SDR & UTMOS & PESQ & ESTOI & SDR & UTMOS \\
\midrule
Unprocessed & 1.97 & 0.79 & \phantom{0}8.4 & 3.11 & 1.88 & 0.60 & $-$4.7 & 2.55 \\
\midrule
MP-SENet~\cite{lu2023mpsenet} & \textbf{3.16} & \textbf{0.88} & \textbf{19.6} & 3.90 & 2.39 & 0.62 & $-$0.3 & 3.00 \\
MetricGAN+~\cite{fu2021metricgan} & 3.13 & 0.83 & \phantom{0}8.6 & 3.63 & 2.53 & 0.62 & $-$5.5 & 2.82 \\
SGMSE+~\cite{richter2023speech} & 2.94 & 0.87 & 17.4 & 3.83 & 2.30 & 0.63 & \phantom{$-$}0.0 & 2.99 \\
SGMSE+ (dereverb)~\cite{richter2023speech} & 1.98 & 0.78 & 10.9 & 3.16 & 1.99 & 0.66 & $-$3.3 & 2.96 \\
NASE~\cite{hu2023nase} & 2.75 & 0.85 & 17.3 & 3.83 & 2.22 & 0.64 & $-$0.5 & 2.99 \\
\midrule
SGMSE+ retrained & 2.85 & 0.86 & 17.1 & 3.92 & 2.60 & 0.77 & \phantom{$-$}2.3 & 3.70 \\
\quad + encoder, input add. & 2.80 & 0.86 & 17.0 & 3.88 & 2.49 & 0.73 & \phantom{$-$}1.4 & 3.62 \\
\quad + encoder, \textbf{deep inj.\ (Proposed)} & 2.83 & 0.86 & 17.4 & \textbf{3.93} & \textbf{2.60} & \textbf{0.80} & \textbf{\phantom{$-$}3.7} & \textbf{3.71} \\
\bottomrule
\end{tabular}
\end{table*}

Table~\ref{tab:main_results} compares our method with baselines.
The three models below the line share the same multi-degradation training data, enabling a controlled comparison of the encoder and injection method.
Compared to unprocessed degraded speech (PESQ 1.88, ESTOI 0.60, SI-SDR $-$4.7\,dB, UTMOS 2.55 on multi-degradation), all enhanced models achieve substantial improvements---with the Proposed model raising PESQ to 2.60 (+0.72), ESTOI to 0.80 (+0.20), and UTMOS to 3.71 (+1.16)---confirming significant enhancement even in challenging compound conditions.

\textbf{Multi-degradation (fair comparison).}
Among models trained on the same data, the injection method is decisive.
Without any encoder (SGMSE+ retrained), the model achieves ESTOI 0.77 and SI-SDR 2.3\,dB.
Adding a WavLM encoder with \textit{input-level addition}---the injection used by NASE~\cite{hu2023nase}---\textit{significantly degrades} performance to ESTOI 0.73 and SI-SDR 1.4\,dB, worse than using no encoder at all ($p < 10^{-7}$, paired $t$-test on 2,472 files).
In contrast, the same encoder with \textit{deep injection} significantly improves to ESTOI 0.80 ($p < 10^{-61}$) and SI-SDR 3.7\,dB ($p < 10^{-15}$).
This demonstrates that deep injection, not the encoder alone, is critical for leveraging degradation information in compound scenarios.

Noise-only baselines (ESTOI 0.62--0.66, SI-SDR $\leq$0) perform substantially worse, and NASE trained on noise-only data achieves only ESTOI 0.64, further confirming that both multi-degradation data and deep injection are necessary.

\textbf{Noise-only.}
On the standard benchmark, specialized models such as MP-SENet (PESQ 3.16) and MetricGAN+ (3.13) achieve higher PESQ, as expected from models designed exclusively for additive noise removal.
All three multi-degradation models achieve comparable noise-only performance (PESQ 2.80--2.85), consistent with the known trade-off between single-task optimization and broader generalization.
Notably, our model produces the highest UTMOS (3.93) across all models, surpassing even MP-SENet (3.90), indicating superior perceptual quality despite the multi-degradation training objective.


\subsection{Ablation Studies}
\label{sec:ablation}

\begin{table}[!t]
\centering
\caption{Ablation study on the multi-degradation test set (2472 files). ``addition'' = input-level addition~\cite{hu2023nase}.}
\label{tab:ablation}
\small
\begin{tabular}{l|cccc}
\toprule
Variant & PESQ & ESTOI & SDR & UTMOS \\
\midrule
\textbf{Proposed (full)} & \textbf{2.60} & \textbf{0.80} & \textbf{3.7} & \textbf{3.71} \\
\quad w/o aux.\ loss ($\lambda\!=\!0$) & 2.58 & 0.77 & 2.2 & 3.63 \\
\quad Input addition & 2.49 & 0.73 & 1.4 & 3.62 \\
No encoder (SGMSE+ retr.) & 2.60 & 0.77 & 2.3 & 3.70 \\
\quad Zero conditioning & 2.14 & 0.67 & $-$1.6 & 3.04 \\
\midrule
\quad Uniform weights & 2.60 & 0.80 & 3.7 & 3.71 \\
\bottomrule
\end{tabular}
\end{table}

Table~\ref{tab:ablation} provides further ablations on the multi-degradation test set.
All variants share the same training data and WavLM encoder (except ``No encoder''), consistent with the controlled comparison in Table~\ref{tab:main_results}.

\textbf{Deep vs.\ shallow injection.}
As also shown in Table~\ref{tab:main_results}, replacing deep injection with input addition~\cite{hu2023nase} degrades every metric: ESTOI $0.80 \to 0.73$, SI-SDR $3.7 \to 1.4$\,dB, UTMOS $3.71 \to 3.62$---\textit{worse than no encoder at all} (ESTOI 0.77, SI-SDR 2.3).
Activation analysis confirms why.
Comparing conditioned ($\mathbf{c}_\text{extra}$ from encoder) and unconditioned ($\mathbf{c}_\text{extra} = \mathbf{0}$) forward passes through all 49 ResBlocks, we measure the relative activation difference $\|\Delta\mathbf{h}_\ell\| / \|\mathbf{h}_\ell\|$ at each layer $\ell$.
Deep injection maintains substantial conditioning influence across the entire network: the mean relative difference is 0.47 averaged over 50 test samples, with 0.43 in the downsampling path (layers 1--15), 0.32 at the bottleneck (16--25), and 0.55 in the upsampling path (26--49)---no monotonic decay.
Input addition, by contrast, perturbs only the input spectrogram; with no per-layer reinforcement, this single perturbation is progressively diluted through residual connections, leaving deeper layers effectively unconditioned.

\textbf{Multi-task loss.}
Removing auxiliary losses ($\lambda = 0$) reduces ESTOI from 0.80 to 0.77 and UTMOS from 3.71 to 3.63.
Per-degradation analysis reveals this effect concentrates on reverberant conditions: for noise+reverb+distortion, ESTOI drops from 0.647 to 0.517 ($-$0.130) and SI-SDR from $-$18.4 to $-$23.8\,dB without auxiliary losses.
This confirms that multi-task supervision strengthens the encoder's reverb discrimination~\cite{caruana1997multitask}.

\textbf{Zero conditioning.}
Setting $\mathbf{c}_\text{extra} = \mathbf{0}$ at inference causes large degradation: PESQ $2.60 \to 2.14$, UTMOS $3.71 \to 3.04$, confirming that the model relies heavily on the conditioning signal.

\textbf{Adaptive vs.\ uniform weighting.}
Scaling each branch by the encoder's predicted confidence versus uniform weights yields identical performance.
This suggests that deep injection already enables the score network to learn which branches to attend to internally, rendering explicit inference-time weighting redundant.


\subsection{Analysis}

\textbf{Encoder head accuracy.}
The multi-task encoder achieves strong degradation characterization: noise binary detection accuracy of 96.7\%, reverb $T_{60}$ correlation of 0.981 (MAE 0.033), and distortion intensity correlation of 0.845 (detection accuracy 86.2\%).
These results confirm that auxiliary losses produce a well-calibrated encoder capable of accurately identifying each degradation type.

\begin{table}[!t]
\centering
\caption{Per-degradation performance breakdown on the multi-degradation test set.}
\label{tab:perdeg}
\small
\setlength{\tabcolsep}{3.5pt}
\begin{tabular}{l|r|cccc}
\toprule
Type & $N$ & PESQ & ESTOI & SDR & UTMOS \\
\midrule
Noise only          & 824 & 2.82 & 0.86 & 17.5  & 3.93 \\
Noise+Reverb        & 330 & 1.72 & 0.65 & $-$17.2 & 3.31 \\
Noise+Distort       & 318 & 2.84 & 0.86 & 15.8  & 3.94 \\
N+R+D               & 338 & 1.74 & 0.65 & $-$18.4 & 3.34 \\
Reverb only         & 341 & 2.01 & 0.74 & $-$17.1 & 3.42 \\
Distort only        & 321 & 4.21 & 0.98 & 23.3  & 4.05 \\
\bottomrule
\end{tabular}
\end{table}

\textbf{Per-degradation breakdown.}
Table~\ref{tab:perdeg} shows performance across degradation types.
Distortion is handled nearly perfectly (PESQ 4.21, UTMOS 4.05), while reverberation causes SI-SDR degradation ($-$17 to $-$18\,dB).
However, UTMOS remains above 3.3 even for reverberant conditions, indicating reasonable perceptual quality despite low SI-SDR.
Single-degradation baselines achieve far worse on reverberant conditions (ESTOI 0.20--0.30, not shown), confirming meaningful enhancement.
The low SI-SDR partly reflects the metric's sensitivity to temporal misalignment~\cite{le2019sdr} rather than actual perceptual quality.

\begin{table}[!t]
\centering
\caption{Evaluation on real-world wild datasets. SGMSE+ (pretrained) is the publicly available denoise model~\cite{richter2023speech}; SGMSE+ (retrained) shares the same multi-degradation training data as Proposed but uses no encoder. SI-SDR is omitted as it is unreliable for reverberant recordings.}
\label{tab:wild}
\small
\begin{tabular}{ll|ccc}
\toprule
Dataset & Model & PESQ & ESTOI & UTMOS \\
\midrule
\multirow{3}{*}{VOiCES} & \textbf{Proposed} & 1.56 & \textbf{0.67} & 2.77 \\
& SGMSE+ (retrained) & \textbf{1.59} & 0.65 & \textbf{2.83} \\
& SGMSE+ (pretrained) & 1.20 & 0.34 & 1.98 \\
\midrule
\multirow{3}{*}{DAPS} & \textbf{Proposed} & \textbf{2.20} & 0.60 & \textbf{3.32} \\
& SGMSE+ (retrained) & 2.04 & 0.65 & 2.96 \\
& SGMSE+ (pretrained) & 2.10 & \textbf{0.69} & 2.96 \\
\midrule
\multirow{3}{*}{URGENT} & \textbf{Proposed} & 1.67 & 0.68 & \textbf{3.10} \\
& SGMSE+ (retrained) & 1.66 & 0.67 & 3.01 \\
& SGMSE+ (pretrained) & \textbf{1.68} & \textbf{0.70} & 2.97 \\
\bottomrule
\end{tabular}
\end{table}

\textbf{Wild dataset generalization.}
Table~\ref{tab:wild} evaluates generalization to real-world recordings with three models: our Proposed, SGMSE+ retrained on the same multi-degradation data (no encoder), and the publicly available SGMSE+ pretrained on noise-only data.
Both multi-degradation models dramatically outperform the pretrained SGMSE+ on VOiCES (PESQ +0.36/+0.39, ESTOI +0.33/+0.31), confirming that diverse training data is the primary driver of wild-data generalization.
While the Proposed model and the retrained baseline yield comparable scores on intrusive metrics (PESQ, ESTOI) across the wild datasets, these reference-based metrics are known to be sub-optimal for heavily mismatched, out-of-domain recordings.
UTMOS---a non-intrusive metric highly correlated with human perception---reveals a clearer picture: the Proposed model achieves the highest UTMOS on DAPS (+0.36) and URGENT (+0.09), demonstrating that the degradation-aware encoder translates to superior perceptual quality even in unseen environments.
Overall, both multi-degradation models substantially outperform the noise-only pretrained baseline in perceptual quality (UTMOS +0.04--0.85).


%==================================================================================
\section{Conclusion}
%==================================================================================

Through controlled experiments where only the injection method varies, we showed that conditioning injection depth critically determines whether degradation-aware conditioning helps or hurts score-based speech enhancement.
Shallow input addition---the standard approach in prior work---performs \textit{worse than no conditioning at all} on compound degradations, while injecting the same conditioning into the timestep embedding yields the best results.
A WavLM encoder with multi-task auxiliary losses produces well-calibrated degradation representations (reverb $T_{60}$ correlation 0.98), enabling a single model to handle noise, reverberation, and distortion jointly.
These results challenge the common assumption that any form of conditioning injection is beneficial, and suggest that \textit{how} conditioning is injected matters as much as \textit{what} is conditioned on---a finding with implications for conditional score-based models beyond speech enhancement.


\section{Generative AI Use Disclosure}
Generative AI tools (Claude) were used for code development assistance and manuscript editing.
All experimental design, analysis, and scientific contributions are the work of the authors.


\bibliographystyle{IEEEtran}
\bibliography{mybib}

\end{document}
